\documentclass{article}

\usepackage{geometry}
\usepackage{indentfirst}
\usepackage{fancyhdr}
\usepackage{graphicx}
\usepackage{titling}

% Automatically prefixes the section number to all first-level enumerate lists
\usepackage{enumitem}
% label=\thesubsection.\arabic* pulls "1.1" and adds the list item number
% leftmargin=3\parindent indents 3 paragraphs.
\setlist[enumerate,1]{label=\thesubsection.\arabic*,leftmargin=\parindent}

\pagestyle{fancy}
\geometry{a4paper, margin=1in}

\begin{document}

\newcommand{\myTitle}{Software Requirement Specification}

\title{\myTitle}
\author{
Project Group 26\\
\\
Jason Reese \\
Juan Torres Llamas \\
Victor Hernandez
}

\date{\today}

\fancyhf{}
\pagestyle{fancy}
\fancyhead[L]{\myTitle}
\fancyhead[R]{Page \thepage}
\fancyfoot[C]{}

\maketitle %
\tableofcontents
\newpage %

\begin{table}[h!]
\centering
\caption{\textbf{Revision History}}
\begin{tabular}{|c|c|c|c|}
\hline
\textbf{Name} & \textbf{Date} & \textbf{Reasons for Changes} & \textbf{Version} \\
\hline
Project Group 26 & 11/17/2026 & First Draft & 1.0.0 \\
\hline

\hline

\hline

\hline
\end{tabular}
\end{table}
\newpage %

\section{Introduction}
\subsection{Purpose}
The Golden Eagle Flight Plan is an application that helps university students track their four-year academic and professional development journey in one centralized platform.
Currently, students must juggle multiple disconnected sources for campus resources and advising: emails, social media, various websites, and one-off meetings with advisors.
This application solves that problem by bringing everything together: a gamified scoreboard system that tracks student progress across Academic Preparation, Career
Development, and Community \& Leadership; a personalized resources and events feed that surfaces content relevant to each student's major, year, and interests; and a faculty-
facing dashboard that allows advisors to post targeted content and monitor individual student progress. The goal is to graduate well-rounded, career-ready, and community-
engaged students by making growth visible and actionable for both students and faculty.

\subsection{Document Conventions}
All requirements are numbered in order to easily establish which requirement is being referred to.

\section{Registration}
	\begin{enumerate}[leftmargin=3\parindent]
	\item When registering, users can choose to sign up as a:
		\begin{itemize}
			\item Student
			\item Faculty
			\item Student Organization
			\item Admin
		\end{itemize}
	\item If the user tries to sign up as a faculty or student organization an admin must review and approve it and assign them roles first.
	\end{enumerate}
\section{Users}
	\subsection{User Roles}
		\begin{enumerate}[leftmargin=3\parindent]
			\item There will be 4 User Roles
			\item Permissions will be granted based on roles.
			\item Student
				\begin{itemize}
					\item The Student role is automatically applied when a user registers as a Student.
				\end{itemize}
			\item Faculty
				\begin{itemize}
					\item The Faculty role requires an Admin to approve the request.
				\end{itemize}
			\item Student Organization
				\begin{itemize}
					\item The Student Organization role requires an Admin to approve the request.
				\end{itemize}
			\item Admin
				\begin{itemize}
					\item The Admin role requires an Admin to approve the request.
					\item When the service is first configured and Admin user must be created.
				\end{itemize}
		\end{enumerate}
	\subsection{User Permissions}
		\begin{enumerate}[leftmargin=3\parindent]
			\item Any registered user should be able to view an events page and RSVP
			\item New users can register to be a student, student organization, or a faculty.
			\item Student organizations and faculty must be approved by an admin.
			\item Only admins and faculty will be able to view a student’s profile.
		\end{enumerate}
		\subsection{Student Users}
			\begin{enumerate}[leftmargin=3\parindent]
				\item Students who have RSVP’d for an event should get a notification when an event they were interested in has been adjusted.
				\item Students will be able to view their own profiles but will not have access to the profiles of other students.
				\item Students will view a homepage upon logging into their accounts.
			\end{enumerate}
		\subsection{Faculty Users}
			\begin{enumerate}[leftmargin=3\parindent]
				\item Faculty and student organizations should be able to create/schedule events, delete/cancel events, reschedule events, and see how many students are interested. 
				\item Faculty should be able to look up a student by name to find their profile.
			\end{enumerate}
		\subsection{Student Organization Users}
			\begin{enumerate}[leftmargin=3\parindent]
				\item Faculty and student organizations should be able to create/schedule events, delete/cancel events, reschedule events, and see how many students are interested.
				\item Student organizations are not allowed to view student profiles due to FERPA.
			\end{enumerate}
		\subsection{Admin Users}
			\begin{enumerate}[leftmargin=3\parindent]
				\item Admin users will be able to manage everything about the site.
				\item Admin users will have all permissions.
			\end{enumerate}
	\subsection{User Profiles}
		\begin{enumerate}[leftmargin=3\parindent]
			\item Every user will have a profile.
			\item Each role will has a specfict profile.
			\item The profile will contain information about the user which is only visible to certain users based on their roles.
			\item Students will only be able to view their own profiles, and faculty or student organization profiles.
		\end{enumerate}
\section{Home Page}
		\subsection{Student Home Page}
			\begin{enumerate}[leftmargin=3\parindent]
				\item The students home page will show events filtered based on the identifiers associated with the student’s profile.
				\item Students will have access to a scoreboard page that will display how much of the flight plan they have completed, suggested activities, and what they should be doing depending on the year they are in.
				\item Students who have RSVP’d for an event should get a notification when an event they were interested in has been adjusted.
			\end{enumerate}
		\subsection{Faculty Home Page}
			\begin{enumerate}[leftmargin=3\parindent]
				\item Faculty 
			\end{enumerate}
		\subsection{Student Organization Home Page}
			\begin{enumerate}[leftmargin=3\parindent]
				\item Student Organization
			\end{enumerate}
		\subsection{Admin Home Page}
			\begin{enumerate}[leftmargin=3\parindent]
				\item The admin 
			\end{enumerate}
\section{Events}
	\begin{enumerate}
		\item When creating events, faculty and student organizations will be able to give the events identifiers.
		\item Students will be able to filter by event identifiers explicitly.
	\end{enumerate}

\section{Other}
	\begin{enumerate}
	\item The flight plan will be split into three sections: academic preparation, career development, and leadership \& community engagement. 
	\item Academic preparation will help the students with figuring out which classes they need to take to graduate and how they can find help with those classes.
	\item Career development will encourage the students to attend events related to their career and will help them find jobs.
	\item Leadership \& Community Engagement will encourage the students to attend events and interact with their community and join clubs or student organizations related to their major.
	\item The design of our project is such that we will have a database that will store our events and our users and their permissions.
	\item We will have a frontend, an API, a backend, and a database.


	\item When creating events, faculty and student organizations will be able to input information related to the event including:
	\item Who is hosting the event
	\item The time the event takes place
	\item If necessary, when to register and when registration opens and closes
	\item Where to go for more information
	\item Whether the event requires registration
	\item The identifiers or keywords associated with the event that can be used to filter for the event.
	\item The location of the event (including online or in-person)
	\item Any additional information regarding the event and who to contact with questions regarding the event.
	\item The ability to upload pictures, flyers or links for the event.
	\item Who posted the event.
	\item Every user will have a password, the password will be salted and use bcrypt.
	\item A future release will support using SingleSignOn instead of storing passwords.
\end{enumerate}



\end{document}
